\documentclass[11pt,a4paper]{article}
\usepackage[T1]{fontenc}
\usepackage{lmodern,array,amsmath,amssymb,booktabs,enumitem,fancyhdr,lastpage,needspace,microtype}
\usepackage[margin=19mm,headheight=14pt,footskip=11mm]{geometry}
\usepackage[hidelinks]{hyperref}
\setlength{\parindent}{0pt}
\setlength{\parskip}{5pt}
\setlength{\emergencystretch}{2em}
\pagestyle{fancy}\fancyhf{}
\fancyhead[L]{\small CS182 --- Lecture 05}
\fancyhead[R]{\small Personal answer copy}
\fancyfoot[L]{\small Independent practice quiz}
\fancyfoot[R]{\small Page \thepage\ of \pageref{LastPage}}
\newcommand{\SetAnswer}[2]{\expandafter\gdef\csname QuizAnswer#1\endcsname{#2}}
\newcommand{\GetAnswer}[1]{\csname QuizAnswer#1\endcsname}
\newlength{\BlankAnswerHeight}\setlength{\BlankAnswerHeight}{16mm}
\newcommand{\AnswerBlank}{\leavevmode\par\vspace{\BlankAnswerHeight}}
% 本文件为独立可填写答题卷；用 pdfLaTeX 编译两次。
% 在 BEGIN TF ANSWERS 和 BEGIN MC ANSWERS 中填写短答案。
% 大题请搜索 BEGIN ANSWER Q25(a) 等注释；每个小问都有独立区域。
% 写完小问后删除该处的 \AnswerBlank，或将 BlankAnswerHeight 改为 0mm。
% 所有打印内容及你填写的回答请使用英文。中文只出现在源文件注释中。
% build_per_lecture.py 不会覆盖已经存在的 Quiz.tex，以保留你的作答。
% BEGIN STUDENT INFO / 个人信息
\newcommand{\StudentName}{}
\newcommand{\StudentID}{}
% END STUDENT INFO
% BEGIN TF ANSWERS / 判断题答题区（第二组花括号填 T 或 F）
\SetAnswer{TF01}{} % Question 1
\SetAnswer{TF02}{} % Question 2
\SetAnswer{TF03}{} % Question 3
\SetAnswer{TF04}{} % Question 4
\SetAnswer{TF05}{} % Question 5
\SetAnswer{TF06}{} % Question 6
\SetAnswer{TF07}{} % Question 7
\SetAnswer{TF08}{} % Question 8
\SetAnswer{TF09}{} % Question 9
\SetAnswer{TF10}{} % Question 10
\SetAnswer{TF11}{} % Question 11
\SetAnswer{TF12}{} % Question 12
% END TF ANSWERS
% BEGIN MC ANSWERS / 不定项选择答题区（填选项字母，如 BD；仅格式示例）
\SetAnswer{MC13}{} % Question 13
\SetAnswer{MC14}{} % Question 14
\SetAnswer{MC15}{} % Question 15
\SetAnswer{MC16}{} % Question 16
\SetAnswer{MC17}{} % Question 17
\SetAnswer{MC18}{} % Question 18
\SetAnswer{MC19}{} % Question 19
\SetAnswer{MC20}{} % Question 20
\SetAnswer{MC21}{} % Question 21
\SetAnswer{MC22}{} % Question 22
\SetAnswer{MC23}{} % Question 23
\SetAnswer{MC24}{} % Question 24
% END MC ANSWERS

\begin{document}
\begin{center}{\Large\bfseries CS182: Introduction to Machine Learning}\par{\Large\bfseries Lecture 05: Bayesian Decision Theory}\par{\large Lecture Quiz and Answer Copy}\end{center}
\textbf{Scope:} Bayesian\_Decision\_Theory.pdf. Independent practice material based on the official 2026 Fall lecture.\par\textbf{Total:} 100 marks. \quad \textbf{Suggested time:} 120 minutes.
\par\textbf{Name:} \underline{\makebox[60mm][l]{\StudentName}}\hfill\textbf{Student ID:} \underline{\makebox[45mm][l]{\StudentID}}
\begin{center}\begin{tabular}{lrrr}\toprule
Section & Questions & Each & Marks\\\midrule
A. Misconceptions: True/False & 12 & 1 & 12\\
B. Multiple response & 12 & 2 & 24\\
C. Proofs and formula derivations & 4 & 10 & 40\\
D. Applications & 2 & 12 & 24\\\midrule
Total & 30 & & 100\\\bottomrule\end{tabular}\end{center}
\textbf{Scoring.} In Section B each item has exactly four options, of which one,
two, or three are correct. The complete correct set earns 2 marks; a nonempty
proper subset earns 1 mark; any incorrect option or no selection earns 0.
T/F answers earn 1 mark only when correct. Show reasoning in Sections C and D;
part marks are available for valid intermediate work. Equivalent correct methods receive credit.
\par\medskip\textbf{Instructions.} Answer in English. Use natural logarithms.
Exact values and justified approximations are accepted. All vectors are columns
unless stated otherwise. Use local assumptions, label conventions, and tie rules.
\par\medskip\textbf{Answer sheet --- Section A}
\begin{center}\renewcommand{\arraystretch}{1.15}\begin{tabular}{|c|c|c|c|c|c|}
\hline
1 & 2 & 3 & 4 & 5 & 6\\\hline
\makebox[20mm]{\rule{0pt}{5mm}\GetAnswer{TF01}} & \makebox[20mm]{\rule{0pt}{5mm}\GetAnswer{TF02}} & \makebox[20mm]{\rule{0pt}{5mm}\GetAnswer{TF03}} & \makebox[20mm]{\rule{0pt}{5mm}\GetAnswer{TF04}} & \makebox[20mm]{\rule{0pt}{5mm}\GetAnswer{TF05}} & \makebox[20mm]{\rule{0pt}{5mm}\GetAnswer{TF06}}\\\hline
\end{tabular}\end{center}
\begin{center}\renewcommand{\arraystretch}{1.15}\begin{tabular}{|c|c|c|c|c|c|}
\hline
7 & 8 & 9 & 10 & 11 & 12\\\hline
\makebox[20mm]{\rule{0pt}{5mm}\GetAnswer{TF07}} & \makebox[20mm]{\rule{0pt}{5mm}\GetAnswer{TF08}} & \makebox[20mm]{\rule{0pt}{5mm}\GetAnswer{TF09}} & \makebox[20mm]{\rule{0pt}{5mm}\GetAnswer{TF10}} & \makebox[20mm]{\rule{0pt}{5mm}\GetAnswer{TF11}} & \makebox[20mm]{\rule{0pt}{5mm}\GetAnswer{TF12}}\\\hline
\end{tabular}\end{center}
\par\medskip\textbf{Answer sheet --- Section B}
\begin{center}\renewcommand{\arraystretch}{1.15}\begin{tabular}{|c|c|c|c|c|c|}
\hline
13 & 14 & 15 & 16 & 17 & 18\\\hline
\makebox[20mm]{\rule{0pt}{5mm}\GetAnswer{MC13}} & \makebox[20mm]{\rule{0pt}{5mm}\GetAnswer{MC14}} & \makebox[20mm]{\rule{0pt}{5mm}\GetAnswer{MC15}} & \makebox[20mm]{\rule{0pt}{5mm}\GetAnswer{MC16}} & \makebox[20mm]{\rule{0pt}{5mm}\GetAnswer{MC17}} & \makebox[20mm]{\rule{0pt}{5mm}\GetAnswer{MC18}}\\\hline
\end{tabular}\end{center}
\begin{center}\renewcommand{\arraystretch}{1.15}\begin{tabular}{|c|c|c|c|c|c|}
\hline
19 & 20 & 21 & 22 & 23 & 24\\\hline
\makebox[20mm]{\rule{0pt}{5mm}\GetAnswer{MC19}} & \makebox[20mm]{\rule{0pt}{5mm}\GetAnswer{MC20}} & \makebox[20mm]{\rule{0pt}{5mm}\GetAnswer{MC21}} & \makebox[20mm]{\rule{0pt}{5mm}\GetAnswer{MC22}} & \makebox[20mm]{\rule{0pt}{5mm}\GetAnswer{MC23}} & \makebox[20mm]{\rule{0pt}{5mm}\GetAnswer{MC24}}\\\hline
\end{tabular}\end{center}
\newpage\section*{Section A. True/False \hfill 12 marks}
\par\noindent\begin{minipage}{\linewidth}
\textbf{1.} A class-conditional density \(p(x\mid C_i)\) must lie between 0 and 1 at every \(x\).
\end{minipage}\par\vspace{3mm}
\par\noindent\begin{minipage}{\linewidth}
\textbf{2.} With unequal misclassification costs, a class whose posterior is below one half can still be the minimum-risk classification action.
\end{minipage}\par\vspace{3mm}
\par\noindent\begin{minipage}{\linewidth}
\textbf{3.} Every point satisfying \(g_i(x)=g_j(x)\) belongs to a decision boundary between \(C_i\) and \(C_j\).
\end{minipage}\par\vspace{3mm}
\par\noindent\begin{minipage}{\linewidth}
\textbf{4.} Replacing every discriminant by \(a(x)g_i(x)+b(x)\), using the same \(a(x)>0\) and \(b(x)\) for all classes, preserves the entire set of maximizing classes.
\end{minipage}\par\vspace{3mm}
\par\noindent\begin{minipage}{\linewidth}
\textbf{5.} Under 0-1 loss, if no features are observed, always predicting a class with the largest prior probability minimizes the probability of error.
\end{minipage}\par\vspace{3mm}
\par\noindent\begin{minipage}{\linewidth}
\textbf{6.} Applying a different strictly increasing function to each class discriminant always preserves the classifier.
\end{minipage}\par\vspace{3mm}
\par\noindent\begin{minipage}{\linewidth}
\textbf{7.} Because the evidence \(p(x)\) cancels when comparing classes, it can also be omitted when reporting normalized posterior probabilities.
\end{minipage}\par\vspace{3mm}
\par\noindent\begin{minipage}{\linewidth}
\textbf{8.} Under 0-1 loss, choosing the largest class likelihood always gives the Bayes decision, even with unequal class priors.
\end{minipage}\par\vspace{3mm}
\par\noindent\begin{minipage}{\linewidth}
\textbf{9.} Randomly mixing classification actions at a fixed \(x\) can give conditional risk strictly below the smallest risk of those actions.
\end{minipage}\par\vspace{3mm}
\par\noindent\begin{minipage}{\linewidth}
\textbf{10.} If two classes tie for the largest posterior probability, choosing either one has the same conditional 0-1 risk.
\end{minipage}\par\vspace{3mm}
\par\noindent\begin{minipage}{\linewidth}
\textbf{11.} With unit error cost and rejection cost \(\lambda\), classification is strictly preferred to rejection exactly when the largest posterior exceeds \(1-\lambda\).
\end{minipage}\par\vspace{3mm}
\par\noindent\begin{minipage}{\linewidth}
\textbf{12.} Naive Bayes requires the features to be independent in the unconditional feature distribution.
\end{minipage}\par\vspace{3mm}
\newpage\section*{Section B. Multiple Response \hfill 24 marks}
\textit{Select all correct options. Each question has 1--3 correct options.}\par
\par\noindent\begin{minipage}{\linewidth}
% BEGIN QUESTION Q13 / 题目与选项区
\textbf{13.} For \(x\in\mathbb R\), define \(g_1(x)=1-x^2\), \(g_2(x)=0\), and \(g_3(x)=1/2\). Use the largest score and allow ties on boundaries. Select every correct statement.
\begin{enumerate}[label=(\Alph*),leftmargin=8mm,itemsep=2pt,topsep=2pt]
\item Class \(C_2\) has a nonempty decision region.
\item The region where \(C_3\) is the unique winner has two disconnected components.
\item The actual class changes occur at \(x=\pm1/\sqrt2\).
\item The equation \(g_1=g_2\) produces \(x=\pm1\), but neither point is a decision boundary.
\end{enumerate}
% END QUESTION Q13

\end{minipage}\par\vspace{3mm}
\par\noindent\begin{minipage}{\linewidth}
% BEGIN QUESTION Q14 / 题目与选项区
\textbf{14.} Let the true joint distribution be known, and use 0-1 loss. Select every correct statement about using features.
\begin{enumerate}[label=(\Alph*),leftmargin=8mm,itemsep=2pt,topsep=2pt]
\item The Bayes error using \(X\) is no larger than the error of always predicting a class with largest prior.
\item Any rule that depends on \(X\) must outperform every constant rule.
\item If each observed input identifies its class with posterior probability 1 almost surely, the Bayes error is zero.
\item If the posterior equals the prior almost surely, observing \(X\) cannot improve the minimum error.
\end{enumerate}
% END QUESTION Q14

\end{minipage}\par\vspace{3mm}
\par\noindent\begin{minipage}{\linewidth}
% BEGIN QUESTION Q15 / 题目与选项区
\textbf{15.} A three-class posterior is \((0.46,0.34,0.20)\). Correct classifications cost 0 and mistakes cost 1. Select every correct statement about a constant rejection cost.
\begin{enumerate}[label=(\Alph*),leftmargin=8mm,itemsep=2pt,topsep=2pt]
\item At rejection cost \(0.55\), acceptance requires the largest posterior to exceed \(0.55\).
\item At rejection cost \(0.50\), rejection is strictly preferable to every class action.
\item At rejection cost \(0.55\), predicting \(C_1\) is strictly preferable to rejection.
\item Increasing the rejection cost cannot enlarge the set where rejection is strictly optimal.
\end{enumerate}
% END QUESTION Q15

\end{minipage}\par\vspace{3mm}
\par\noindent\begin{minipage}{\linewidth}
% BEGIN QUESTION Q16 / 题目与选项区
\textbf{16.} Assume all discriminants are finite. Which transformations preserve the full maximizing-class set at every \(x\), for arbitrary original scores?
\begin{enumerate}[label=(\Alph*),leftmargin=8mm,itemsep=2pt,topsep=2pt]
\item Multiply every score by the same negative constant.
\item Add a different constant \(b_i\) to each class score.
\item Apply one common strictly increasing function \(f\) to every score.
\item Replace every \(g_i(x)\) by \(a(x)g_i(x)+b(x)\), with common \(a(x)>0\) and \(b(x)\).
\end{enumerate}
% END QUESTION Q16

\end{minipage}\par\vspace{3mm}
\newpage\section*{Section B. Multiple Response \hfill 24 marks}
\textit{Select all correct options. Each question has 1--3 correct options.}\par
\par\noindent\begin{minipage}{\linewidth}
% BEGIN QUESTION Q17 / 题目与选项区
\textbf{17.} Consider any finite loss matrix \(\lambda_{ik}\). Which modifications preserve all minimum-risk actions for every posterior?
\begin{enumerate}[label=(\Alph*),leftmargin=8mm,itemsep=2pt,topsep=2pt]
\item Add one state-dependent number \(b_k\) to every entry in column \(k\).
\item Replace risks by their negatives and choose the largest resulting score.
\item Multiply the entire loss matrix by one positive constant.
\item Add a different arbitrary constant to every row.
\end{enumerate}
% END QUESTION Q17

\end{minipage}\par\vspace{3mm}
\par\noindent\begin{minipage}{\linewidth}
% BEGIN QUESTION Q18 / 题目与选项区
\textbf{18.} For two classes with positive priors and positive likelihoods, let \(P(C_1)=1/4\), \(P(C_2)=3/4\), and \(L(x)=p(x\mid C_1)/p(x\mid C_2)\). Under 0-1 loss, select every correct statement.
\begin{enumerate}[label=(\Alph*),leftmargin=8mm,itemsep=2pt,topsep=2pt]
\item For the given priors, a log discriminant is \(\log L(x)-\log3\).
\item Predicting \(C_1\) is strictly optimal when \(L(x)>3\).
\item The evidence is needed to determine which posterior is larger.
\item Replacing the priors by \(0.1\) and \(0.9\) changes the likelihood-ratio threshold to 9.
\end{enumerate}
% END QUESTION Q18

\end{minipage}\par\vspace{3mm}
\par\noindent\begin{minipage}{\linewidth}
% BEGIN QUESTION Q19 / 题目与选项区
\textbf{19.} Which statements about the naive Bayes assumption are correct?
\begin{enumerate}[label=(\Alph*),leftmargin=8mm,itemsep=2pt,topsep=2pt]
\item It is sufficient that the features be mutually independent conditional on each class.
\item Unconditional independence alone guarantees the required class-conditional factorization.
\item The class prior is unnecessary once the likelihood is factorized.
\item The factorization makes \(P(C_i\mid x)\) equal to \(\prod_jP(C_i\mid x_j)\) without any normalization or prior correction.
\end{enumerate}
% END QUESTION Q19

\end{minipage}\par\vspace{3mm}
\par\noindent\begin{minipage}{\linewidth}
% BEGIN QUESTION Q20 / 题目与选项区
\textbf{20.} For three classes, let the priors be \((0.2,0.3,0.5)\) and the class likelihoods at one observed \(x\) be \((0.9,0.4,0.2)\). Under 0-1 loss, select every correct statement.
\begin{enumerate}[label=(\Alph*),leftmargin=8mm,itemsep=2pt,topsep=2pt]
\item The Bayes action predicts \(C_1\).
\item The minimum conditional probability of error is \(0.45\).
\item The posterior probability of \(C_1\) is \(0.60\).
\item The evidence is \(p(x)=0.40\).
\end{enumerate}
% END QUESTION Q20

\end{minipage}\par\vspace{3mm}
\newpage\section*{Section B. Multiple Response \hfill 24 marks}
\textit{Select all correct options. Each question has 1--3 correct options.}\par
\par\noindent\begin{minipage}{\linewidth}
% BEGIN QUESTION Q21 / 题目与选项区
\textbf{21.} A posterior is \((0.4,0.4,0.2)\), unit error loss is used, and rejection costs \(0.6\). Select every correct statement.
\begin{enumerate}[label=(\Alph*),leftmargin=8mm,itemsep=2pt,topsep=2pt]
\item Any random mixture of rejection and predictions of \(C_1\) or \(C_2\) also has risk \(0.6\).
\item Predicting \(C_3\) is also optimal because all posterior values are below \(0.6\).
\item The minimum conditional risk is \(0.6\).
\item Rejection is the unique risk-minimizing action.
\end{enumerate}
% END QUESTION Q21

\end{minipage}\par\vspace{3mm}
\par\noindent\begin{minipage}{\linewidth}
% BEGIN QUESTION Q22 / 题目与选项区
\textbf{22.} Write \(q=P(C_1\mid x)=0.55\). Correct decisions cost 0; predicting \(C_1\) when \(C_2\) is true costs 8, and predicting \(C_2\) when \(C_1\) is true costs 2. Select every correct statement.
\begin{enumerate}[label=(\Alph*),leftmargin=8mm,itemsep=2pt,topsep=2pt]
\item Predicting \(C_1\) has risk \(3.6\).
\item Predicting \(C_2\) minimizes conditional risk.
\item Without rejection, predicting \(C_1\) is strictly optimal when \(q>0.8\).
\item The larger posterior must determine the minimum-risk action.
\end{enumerate}
% END QUESTION Q22

\end{minipage}\par\vspace{3mm}
\par\noindent\begin{minipage}{\linewidth}
% BEGIN QUESTION Q23 / 题目与选项区
\textbf{23.} At a fixed input, suppose the risks of three available actions are \(0.2,0.5,0.2\). Select every correct statement about randomized actions.
\begin{enumerate}[label=(\Alph*),leftmargin=8mm,itemsep=2pt,topsep=2pt]
\item Assigning positive probability to the second action can still yield risk below \(0.2\).
\item Every minimum-risk randomization must put zero probability on the second action.
\item A half-and-half mixture of the first and third actions has risk \(0.2\).
\item The existence of two minimizing actions forces a unique optimal randomization.
\end{enumerate}
% END QUESTION Q23

\end{minipage}\par\vspace{3mm}
\par\noindent\begin{minipage}{\linewidth}
% BEGIN QUESTION Q24 / 题目与选项区
\textbf{24.} At an observed point, two continuous class densities have values \(4\) and \(1\), and their priors are \(0.1\) and \(0.9\). Under 0-1 loss, select every correct statement.
\begin{enumerate}[label=(\Alph*),leftmargin=8mm,itemsep=2pt,topsep=2pt]
\item The evidence density is \(1.3\).
\item The posterior of \(C_1\) is \(4/13\).
\item The Bayes action predicts \(C_2\).
\item The larger likelihood proves that \(C_1\) is the Bayes choice.
\end{enumerate}
% END QUESTION Q24

\end{minipage}\par\vspace{3mm}
\newpage\section*{Section C. Proofs and Formula Derivations}
\textbf{25. Why pointwise decisions minimize total risk}\hfill\textbf{[10 marks]}\par\medskip
\begin{enumerate}[label=(\alph*),leftmargin=8mm,itemsep=3pt]
\item Let \(A\) be a finite action set, let \(\lambda_{ak}\) be finite nonnegative losses, and let \(X\) have density or mass function \(p(x)\). Define a deterministic rule \(\delta(x)\). Starting from \(E[\lambda_{\delta(X),C}]\), derive the total-risk formula and prove that choosing a minimum-conditional-risk action at each \(x\) minimizes total risk. State the almost-sure qualification.\hfill\textbf{[4 marks]}
\item Specialize to \(K\) classification actions and 0-1 loss. Derive the MAP rule and the Bayes error expression in terms of the largest posterior.\hfill\textbf{[3 marks]}
\item Prove that this Bayes error is at most \(1-\max_iP(C_i)\), the best constant-rule error. Explain why this does not guarantee that every fitted classifier improves on the constant rule.\hfill\textbf{[3 marks]}
\end{enumerate}\par\medskip\hrule\medskip\textbf{Student response}\par
% BEGIN ANSWER Q25 / 第25题答题区
\begin{enumerate}[label=(\alph*),leftmargin=8mm,itemsep=3pt]
\item
% >>> BEGIN ANSWER Q25(a)
% 在这里写英文推导或证明；完成后可删除下方空白占位命令。
\AnswerBlank
% <<< END ANSWER Q25(a)
\item
% >>> BEGIN ANSWER Q25(b)
% 在这里写英文推导或证明；完成后可删除下方空白占位命令。
\AnswerBlank
% <<< END ANSWER Q25(b)
\item
% >>> BEGIN ANSWER Q25(c)
% 在这里写英文推导或证明；完成后可删除下方空白占位命令。
\AnswerBlank
% <<< END ANSWER Q25(c)
\end{enumerate}
% END ANSWER Q25
\newpage\section*{Section C. Proofs and Formula Derivations}
\textbf{26. Asymmetric errors and rejection thresholds}\hfill\textbf{[10 marks]}\par\medskip
\begin{enumerate}[label=(\alph*),leftmargin=8mm,itemsep=3pt]
\item For two classes, write \(q=P(C_1\mid x)\). Correct decisions cost 0, predicting \(C_1\) when \(C_2\) is true costs \(a>0\), and predicting \(C_2\) when \(C_1\) is true costs \(b>0\). Derive the optimal posterior threshold without rejection and the equivalent likelihood-ratio threshold when priors and likelihoods are positive.\hfill\textbf{[4 marks]}
\item Add a reject action with constant loss \(r>0\) for either true class. Derive the interval where rejection is strictly optimal and determine when that interval has positive width.\hfill\textbf{[3 marks]}
\item Return to \(K\) classes, unit cost for any error, and reject cost \(\lambda\). Derive the acceptance rule. Then show that if \(\lambda>1-1/K\), rejection cannot be strictly optimal at any input.\hfill\textbf{[3 marks]}
\end{enumerate}\par\medskip\hrule\medskip\textbf{Student response}\par
% BEGIN ANSWER Q26 / 第26题答题区
\begin{enumerate}[label=(\alph*),leftmargin=8mm,itemsep=3pt]
\item
% >>> BEGIN ANSWER Q26(a)
% 在这里写英文推导或证明；完成后可删除下方空白占位命令。
\AnswerBlank
% <<< END ANSWER Q26(a)
\item
% >>> BEGIN ANSWER Q26(b)
% 在这里写英文推导或证明；完成后可删除下方空白占位命令。
\AnswerBlank
% <<< END ANSWER Q26(b)
\item
% >>> BEGIN ANSWER Q26(c)
% 在这里写英文推导或证明；完成后可删除下方空白占位命令。
\AnswerBlank
% <<< END ANSWER Q26(c)
\end{enumerate}
% END ANSWER Q26
\newpage\section*{Section C. Proofs and Formula Derivations}
\textbf{27. Equivalent discriminants and the role of strict monotonicity}\hfill\textbf{[10 marks]}\par\medskip
\begin{enumerate}[label=(\alph*),leftmargin=8mm,itemsep=3pt]
\item Prove that applying one common strictly increasing function to all finite discriminants preserves their full argmax set. Explain why merely nondecreasing is insufficient for preserving the argmax set, using a concrete counterexample.\hfill\textbf{[4 marks]}
\item For \(p(x)>0\) and positive class joint scores, derive the equivalence of posterior, prior-weighted likelihood, and log prior-weighted likelihood discriminants. Identify the qualification needed for zero joint scores.\hfill\textbf{[3 marks]}
\item For two classes with positive priors and likelihoods, derive a single scalar log discriminant. Show how its sign determines the 0-1 decision, and explain why a negative-risk discriminant rather than posterior probability is needed for arbitrary losses.\hfill\textbf{[3 marks]}
\end{enumerate}\par\medskip\hrule\medskip\textbf{Student response}\par
% BEGIN ANSWER Q27 / 第27题答题区
\begin{enumerate}[label=(\alph*),leftmargin=8mm,itemsep=3pt]
\item
% >>> BEGIN ANSWER Q27(a)
% 在这里写英文推导或证明；完成后可删除下方空白占位命令。
\AnswerBlank
% <<< END ANSWER Q27(a)
\item
% >>> BEGIN ANSWER Q27(b)
% 在这里写英文推导或证明；完成后可删除下方空白占位命令。
\AnswerBlank
% <<< END ANSWER Q27(b)
\item
% >>> BEGIN ANSWER Q27(c)
% 在这里写英文推导或证明；完成后可删除下方空白占位命令。
\AnswerBlank
% <<< END ANSWER Q27(c)
\end{enumerate}
% END ANSWER Q27
\newpage\section*{Section C. Proofs and Formula Derivations}
\textbf{28. What a decision boundary actually represents}\hfill\textbf{[10 marks]}\par\medskip
\begin{enumerate}[label=(\alph*),leftmargin=8mm,itemsep=3pt]
\item Suppose \(g_i(x)=w_i^\mathsf Tx+b_i\) on \(\mathbb R^d\). Show that the weak decision region \(\mathcal R_i=\{x:g_i(x)\ge g_j(x)\text{ for all }j\}\) is convex. Describe the supporting equality between two competing classes.\hfill\textbf{[4 marks]}
\item Give an explicit two-class example on \(\mathbb R\) with a disconnected decision region, and derive the regions and boundaries. Use \(g_1(x)=x^2-4\) and \(g_2(x)=0\).\hfill\textbf{[3 marks]}
\item For \(g_1(x)=2-x^2\), \(g_2(x)=0\), and \(g_3(x)=1\), identify all pairwise equality points and determine which are actual decision boundaries. Use the result to disprove the claim that every discriminant equality is a boundary.\hfill\textbf{[3 marks]}
\end{enumerate}\par\medskip\hrule\medskip\textbf{Student response}\par
% BEGIN ANSWER Q28 / 第28题答题区
\begin{enumerate}[label=(\alph*),leftmargin=8mm,itemsep=3pt]
\item
% >>> BEGIN ANSWER Q28(a)
% 在这里写英文推导或证明；完成后可删除下方空白占位命令。
\AnswerBlank
% <<< END ANSWER Q28(a)
\item
% >>> BEGIN ANSWER Q28(b)
% 在这里写英文推导或证明；完成后可删除下方空白占位命令。
\AnswerBlank
% <<< END ANSWER Q28(b)
\item
% >>> BEGIN ANSWER Q28(c)
% 在这里写英文推导或证明；完成后可删除下方空白占位命令。
\AnswerBlank
% <<< END ANSWER Q28(c)
\end{enumerate}
% END ANSWER Q28
\newpage\section*{Section D. Applications}
\textbf{29. A sensor with asymmetric mistakes and a review action}\hfill\textbf{[12 marks]}\par\medskip
\begin{enumerate}[label=(\alph*),leftmargin=8mm,itemsep=3pt]
\item A sensor reports \(A,B,\) or \(C\). The class priors are \(P(C_1)=0.35\), \(P(C_2)=0.65\), and the class-conditional probabilities for \((A,B,C)\) are \((0.50,0.30,0.20)\) in \(C_1\) and \((0.10,0.40,0.50)\) in \(C_2\). Compute each report probability and each posterior \(q_s=P(C_1\mid s)\).\hfill\textbf{[4 marks]}
\item Correct classifications cost 0. Predicting \(C_1\) when \(C_2\) is true costs 3; predicting \(C_2\) when \(C_1\) is true costs 9. A human review action costs 1.7 regardless of class. Compute all three conditional risks for each report and choose the optimal action.\hfill\textbf{[4 marks]}
\item Compute the expected loss of your review-enabled rule. Compare it with the best rule forced to classify every report and with the MAP rule that ignores the asymmetric costs.\hfill\textbf{[4 marks]}
\end{enumerate}\par\medskip\hrule\medskip\textbf{Student response}\par
% BEGIN ANSWER Q29 / 第29题答题区
\begin{enumerate}[label=(\alph*),leftmargin=8mm,itemsep=3pt]
\item
% >>> BEGIN ANSWER Q29(a)
% 在这里写英文推导或证明；完成后可删除下方空白占位命令。
\AnswerBlank
% <<< END ANSWER Q29(a)
\item
% >>> BEGIN ANSWER Q29(b)
% 在这里写英文推导或证明；完成后可删除下方空白占位命令。
\AnswerBlank
% <<< END ANSWER Q29(b)
\item
% >>> BEGIN ANSWER Q29(c)
% 在这里写英文推导或证明；完成后可删除下方空白占位命令。
\AnswerBlank
% <<< END ANSWER Q29(c)
\end{enumerate}
% END ANSWER Q29
\newpage\section*{Section D. Applications}
\textbf{30. Conditional independence, posterior decisions, and rejection}\hfill\textbf{[12 marks]}\par\medskip
\begin{enumerate}[label=(\alph*),leftmargin=8mm,itemsep=3pt]
\item Two binary features are independent conditional on class. Let \(P(C_1)=0.4\), \(P(C_2)=0.6\); within \(C_1\), the probabilities \(P(X_1=1\mid C_1)\), \(P(X_2=1\mid C_1)\) are \(0.8,0.3\); within \(C_2\), they are \(0.25,0.7\). For every binary pair \((00,01,10,11)\), compute the two joint class-feature scores, give the MAP class, and find the overall 0-1 Bayes error.\hfill\textbf{[4 marks]}
\item Add a rejection action costing \(0.30\), with unit misclassification costs. Determine which feature pairs are accepted and compute the total expected risk including rejected inputs.\hfill\textbf{[4 marks]}
\item Are \(X_1\) and \(X_2\) independent without conditioning on class? Establish your answer by computing the necessary marginal and joint probabilities.\hfill\textbf{[4 marks]}
\end{enumerate}\par\medskip\hrule\medskip\textbf{Student response}\par
% BEGIN ANSWER Q30 / 第30题答题区
\begin{enumerate}[label=(\alph*),leftmargin=8mm,itemsep=3pt]
\item
% >>> BEGIN ANSWER Q30(a)
% 在这里写英文推导或证明；完成后可删除下方空白占位命令。
\AnswerBlank
% <<< END ANSWER Q30(a)
\item
% >>> BEGIN ANSWER Q30(b)
% 在这里写英文推导或证明；完成后可删除下方空白占位命令。
\AnswerBlank
% <<< END ANSWER Q30(b)
\item
% >>> BEGIN ANSWER Q30(c)
% 在这里写英文推导或证明；完成后可删除下方空白占位命令。
\AnswerBlank
% <<< END ANSWER Q30(c)
\end{enumerate}
% END ANSWER Q30
\end{document}
